Os ciclones extratropicais (ETCs) geram ventos intensos e precipitações volumosas na costa do Atlântico Sudoccidental, cuja coincidência espacial costuma ser assumida sem verificação explícita. Enquanto o Hemisfério Norte conta com estudos sistemáticos sobre como esses campos se distribuem, o Hemisfério Sul carece de uma integração equivalente que combine densidade espacial e análise de Funções Ortogonais Empíricas (EOF) univariadas e multivariadas por fase do ciclo de vida. Esta tese desenvolve tal integração para avaliar se a intensidade dinâmica do ETC prediz a magnitude e a localização de seus máximos de vento a 10 metros e precipitação, ou se ambos os campos desenvolvem uma organização estrutural qualitativamente distinta segundo a fase evolutiva.

A análise se constrói sobre o catálogo de trajetórias por vorticidade relativa a 850~hPa \citep{gramcianinov2020analysis} e a classificação objetiva do ciclo de vida em quatro fases mediante \textit{CycloPhaser} \citep{coutodesouza2024new, deSouza2025CycloPhaser}, sobre o ERA5 (1979--2020). A amostra foi estratificada por três regiões de ciclogênese (SBR, LPB, ARG) e dois estratos de intensidade: a População Global e o subconjunto dos ciclones mais intensos (p90, extraído da População Global). Sobre essa base, foram aplicadas, de forma progressiva, funções de densidade de probabilidade (PDF) de magnitude, estimação de densidade espacial (PDFe) de localização, decomposição EOF univariada com avaliação de significância mediante Bootstrap-$t$, e decomposição MEOF da covariância conjunta, integradas finalmente em prototipos estruturais.

Na População Global, o vento se comporta de forma passiva e previsível, com um único modo EOF dominante (31--57\% da variância) correlacionado com a vorticidade; a precipitação, por sua vez, se distribui entre múltiplos modos (EOF1 apenas 11--16\%). Essa arquitetura se reorganiza drasticamente no subconjunto p90: o vento se concentra fortemente no quadrante noroeste (moda de 23--25~m\,s$^{-1}$, densidades de até $27\times10^{-4}$ em SBR), enquanto a precipitação colapsa no núcleo e se reorganiza em uma estrutura em arco em direção ao sudeste e sul. O MEOF mostra que ambas as transformações são uma mesma reorganização: na População Global captura 14.8--26.6\% da covariância conjunta, com os gradientes de precipitação e vento orientados em eixos perpendiculares, mas em p90 cai para 11.65--23.47\% e revela, a partir da maturidade, um padrão de anticovariância geométrica, respaldado por uma correlação significativa entre as componentes principais de ambos os campos ($|r|$ entre $0.31$ e $0.45$). Os prototipos estruturais quantificam essa separação em mais de $8^{\circ}$--$9^{\circ}$ durante a maturidade e o decaimento, com ARG alcançando a maior magnitude (${\sim}10$--$11^{\circ}$). A magnitude dessa separação varia regionalmente sem um único fator identificado, embora sua direção, precipitação ao sudeste, vento ao noroeste, se mantenha consistente nas três regiões.

Esses resultados mostram que, nos ETCs maduros do Atlântico Sudoccidental, os máximos de vento e precipitação operam em quadrantes opostos, e que sua parametrização mediante métricas escalares de intensidade central subestima sistematicamente tanto a magnitude quanto a localização de ambos os forçantes. Este trabalho combina PDFe espacial, EOF univariado e MEOF para quantificar essa separação por fase evolutiva e região de gênese no Hemisfério Sul, com implicações para a gestão do risco costeiro no Brasil, Uruguai e Argentina.

\vspace{1cm}

\noindent \textbf{Palavras-chave:} South America; extratropical cyclones; wind at 10 meters; precipitation; spatial structure; multivariate EOF.
